%%%%%%%%%%%%%%%%%%%%%%%%%%%%%%%%%%%%%%%%%
% a0poster Portrait Poster
% LaTeX Template
% Version 1.0 (22/06/13)
%
% The a0poster class was created by:
% Gerlinde Kettl and Matthias Weiser (tex@kettl.de)
%
% This template has been downloaded from:
% http://www.LaTeXTemplates.com
%
% License:
% CC BY-NC-SA 3.0 (http://creativecommons.org/licenses/by-nc-sa/3.0/)
%
%%%%%%%%%%%%%%%%%%%%%%%%%%%%%%%%%%%%%%%%%

%----------------------------------------------------------------------------------------
%	PACKAGES AND OTHER DOCUMENT CONFIGURATIONS
%----------------------------------------------------------------------------------------

\documentclass[a0,portrait,15pt]{a0poster}

\usepackage{multicol} % This is so we can have multiple columns of text side-by-side
\columnsep=100pt % This is the amount of white space between the columns in the poster
\columnseprule=3pt % This is the thickness of the black line between the columns in the poster

\usepackage[svgnames]{xcolor} % Specify colors by their 'svgnames', for a full list of all colors available see here: http://www.latextemplates.com/svgnames-colors

\usepackage{times} % Use the times font
%\usepackage{palatino} % Uncomment to use the Palatino font

\usepackage{graphicx} % Required for including images
\graphicspath{{figures/}} % Location of the graphics files
\usepackage{booktabs} % Top and bottom rules for table
\usepackage[font=small,labelfont=bf]{caption} % Required for specifying captions to tables and figures

\DeclareCaptionFont{myblue}{\color{RoyalBlue}}
\captionsetup{labelfont={myblue,bf}}

\usepackage{amsfonts, amsmath, amsthm, amssymb} % For math fonts, symbols and environments
\usepackage{wrapfig} % Allows wrapping text around tables and figures
\usepackage[position=top]{subfig}
\graphicspath{{figures/}{../paper/figures/}{./}}
\usepackage{tikz,pgf}
\usetikzlibrary{arrows,positioning}

\begin{document}

%----------------------------------------------------------------------------------------
%	POSTER HEADER
%----------------------------------------------------------------------------------------

% The header is divided into two boxes:
% The first is 75% wide and houses the title, subtitle, names, university/organization and contact information
% The second is 25% wide and houses a logo for your university/organization or a photo of you
% The widths of these boxes can be easily edited to accommodate your content as you see fit


\begin{minipage}[b]{0.75\linewidth}
\veryHuge \color{NavyBlue} \textbf{Bandwidth extension of musical audio signals with no side information using dilated convolutional neural networks} \color{Black}\\ % Title
\Huge\textit{}\\[2cm] % Subtitle
\huge \textbf{Mathieu Lagrange, F\'elix Gontier}\\[0.5cm] % Author(s)
\huge LS2N, CNRS, \'Ecole Centrale Nantes\\[0.4cm] % University/organization
\end{minipage}
%
\begin{minipage}[b]{0.25\linewidth}
\includegraphics[width=20cm]{figures/logoLs2n.jpg}\\
\end{minipage}

\vspace{1cm} % A bit of extra whitespace between the header and poster content

%----------------------------------------------------------------------------------------

\begin{multicols}{2} % This is how many columns your poster will be broken into, a portrait poster is generally split into 2 columns

%----------------------------------------------------------------------------------------
%	ABSTRACT
%----------------------------------------------------------------------------------------

\color{Navy} % Navy color for the abstract

\begin{abstract}

Bandwidth extension has a long history in audio processing. While speech processing tools do not rely on side information, production-ready bandwidth extension tools of general audio signals rely on side information that has to be transmitted alongside the bitstream of the low frequency part, mostly because polyphonic music has a more complex and less predictable spectral structure than speech.

This paper studies the benefit of considering \textbf{a dilated fully convolutional neural network} to perform the bandwidth extension of musical audio signals \textbf{with no side information} on the \textbf{magnitude spectra}. Experimental evaluation using two public datasets, \textit{medley-solos-db} and \textit{gtzan}, respectively of monophonic and polyphonic music demonstrate that the proposed architecture achieves state of the art performance.

\end{abstract}

%----------------------------------------------------------------------------------------
%	INTRODUCTION
%----------------------------------------------------------------------------------------

\color{SaddleBrown} % SaddleBrown color for the introduction

\section*{Background}
\nocite{*} % Print all references regardless of whether they were cited in the poster or not
\bibliographystyle{plain} % Plain referencing style
\bibliography{bib} % Use the example bibliography file sample.bib


%----------------------------------------------------------------------------------------
%	OBJECTIVES
%----------------------------------------------------------------------------------------
\color{DarkSlateGray}


\color{SaddleBrown}


\color{DarkSlateGray} % DarkSlateGray color for the rest of the

\section*{}


\section*{Results}



\color{SaddleBrown} % SaddleBrown color for the conclusions to make them stand out

\section*{Conclusions}



\color{DarkSlateGray} % Set the color back to DarkSlateGray for the rest of the content

%----------------------------------------------------------------------------------------
%	FORTHCOMING RESEARCH
%----------------------------------------------------------------------------------------

\section*{Forthcoming Research}

 %----------------------------------------------------------------------------------------
%	REFERENCES
%----------------------------------------------------------------------------------------


%----------------------------------------------------------------------------------------
%	ACKNOWLEDGEMENTS
%----------------------------------------------------------------------------------------

\section*{Acknowledgements}


The authors would like to acknowledge support for this project
from ANR project Houle (grant ANR-11-JS03-005-01) and ANR project Cense (grant ANR-16-CE22-0012).

%----------------------------------------------------------------------------------------

\end{multicols}
\end{document}
